\documentclass[11pt]{article}
\usepackage[a4paper,margin=21mm]{geometry}
\usepackage{fontspec}
\setmainfont{Latin Modern Roman}
\newfontfamily\CJKfont{Noto Serif CJK SC}
\XeTeXlinebreaklocale "zh"
\XeTeXlinebreakskip=0pt plus 1pt
\usepackage{amsmath,amssymb,amsthm,amscd,graphicx,color}
\usepackage{xurl}
\usepackage[unicode,hidelinks]{hyperref}
\allowdisplaybreaks
\setlength\emergencystretch{3em}
\numberwithin{equation}{section}

\numberwithin{equation}{section}

\newtheorem{Theorem}{Theorem}[section]
\newtheorem{Corollary}[Theorem]{Corollary}
\newtheorem{Lemma}[Theorem]{Lemma}
\newtheorem{Proposition}[Theorem]{Proposition}
\newtheorem{Fact}[Theorem]{Fact}
 { \theoremstyle{definition}
\newtheorem{Definition}[Theorem]{Definition}
\newtheorem{Note}[Theorem]{Note}
\newtheorem{Example}[Theorem]{Example}
\newtheorem{Remark}[Theorem]{Remark}
\newtheorem{Setting}[Theorem]{Setting}
\newtheorem{Notation}[Theorem]{Notation}
}


\makeatletter\newlabel{mthm1}{{1.1}{}{Original source locator}{}{}}\makeatother
\makeatletter\newlabel{proper}{{4.5}{}{Original source locator}{}{}}\makeatother
\begin{document}
\begin{center}\Large L2 eigenvalue multiplicities along lambda\_m grow without bound on discontinuous AdS3 quotients satisfying the uniform exponential orbit-growth condition\end{center}
\noindent\textbf{Stable ID:} 1409\quad\textbf{Author:} Kazuki Kannaka\par
\noindent\textbf{Work:} Linear Independence of Generalized Poincaré Series for Anti-de Sitter 3-Manifolds\par
\noindent\textbf{Source:} \url{https://sigma-journal.com/2021/042/}\par
\noindent\textbf{Version:} Official SIGMA 17 (2021) article 042, DOI 10.3842/SIGMA.2021.042; journal PDF and native TeX acquired 2026-09-30, exact bytes pinned by SHA256\par
\noindent\textbf{Rights:} CC BY 4.0. Authors retain copyright. \url{https://creativecommons.org/licenses/by/4.0/}. Reuse and typesetting adaptation; original language and mathematical content preserved.\par
\noindent\textbf{Difficulty (prior selection preserved):} {\CJKfont H2: The proof constructs arbitrarily many independent generalized Poincare series by arranging their angular phases with a lacunary ternary sign pattern. A coefficient-positive polynomial then permits a uniform bound on all nonidentity orbit terms, strictly smaller than the selected identity contribution. This handles linear combinations of functions vanishing at the origin and yields multiplicity growth beyond the earlier nonvanishing result for a single series.}\par
\medskip\noindent\textbf{Editorial boundary note (not author text):} Selected general Theorem3.1 for a free properly discontinuous AdS3 quotient satisfying uniform exponential orbit growth2.9 at every point/radius, with lambda\_m=4m(m-1). Full weak-L2/Lorentzian/quadratic/coordinate/pseudo-distance/spherical-function definitions and exact established Kassel–Kobayashi convergence fact2.4 are retained, followed by complete written separation/Baire argument, new lacunary phase lemma/positive polynomial decay lemma/explicit m\_Gamma(k), all nonidentity-tail and identity-term comparisons and final independence/multiplicity proof. This is divergence of multiplicity along lambda\_m, not infinite multiplicity at each m or an elliptic Weyl law. Original label/type/index variants, including the native Proposition pointer to the displayed threshold label, are unchanged. Finitely-generated, stable-deformation/orbifold consequences are unselected; no supporting lemma receives an extra count. Prior convergence and hyperbolic background remain exact references, no assistant construction.\par
\noindent\textbf{External original reference mthm1:} Original introductory finitely-generated corollary Theorem1.1/native117–127; selected general exponential-orbit-growth Theorem3.1 and its full proof are retained.\par
\noindent\textbf{External original reference proper:} Original Fact4.5/native938–949 quotes established sharpness of finitely-generated groups; it is mentioned only in the unselected finitely-generated contextual consequence, not an extra premise of selected3.1.\par
\medskip\hrule\medskip\noindent\textbf{Author text begins below.}\par

% Original source lines 53--115
\section{Introduction}
A pseudo-Riemannian manifold is a smooth manifold $M$ equipped with
a smooth non-degenerate symmetric bilinear tensor $g$
 of signature $(p,q)$ on $M$.
It is called Riemannian if $q=0$, and Lorentzian if $q=1$.
As in the Riemannian case,
the Laplacian $\square_{M}:=\operatorname{div}_{M} \circ \operatorname{grad}_{M}$
is defined as a second-order differential operator on $M$. We~note that it is a
hyperbolic differential operator if $M$ is Lorentzian. We~write $L^2(M)$ for the Hilbert space of square-integrable functions on $M$
with respect to the Radon measure induced by the pseudo-Riemannian structure.
For $\lambda \in \mathbb{C}$, we denote by
\begin{gather*}
L^2_{\lambda}(M) :=
\big\{f \in L^2(M) \mid \square_{M}f = \lambda f \text{ in the weak sense}\big\}.
\end{gather*}
The set of $L^2$-eigenvalues
$\mathrm{Spec}_{d}(\square_{M}) := \big\{ \lambda \in \mathbb{C}
\mid L^2_{\lambda}(M)\neq0\big\}$
is called the \textit{discrete spectrum} of~$\square_{M}$.

Our interest is the multiplicities of $L^2$-eigenvalues $\lambda$ of $\square_{M}$,
denoted by
\begin{gather*}
\mathcal{N}_{M}(\lambda):= \operatorname{dim}_{\mathbb{C}}L^2_{\lambda}(M)
\in \mathbb{N}\cup\{\infty\}.
\end{gather*}
In the Riemannian case, the Laplacian is an elliptic differential operator
and the distribution of~its discrete spectrum has been investigated extensively,
such as the Weyl law for compact Riemannian manifolds.
However, it is not the case for non-Riemannian manifolds.
Kobayashi~\cite{Kob16}, and later Fox--Strichartz~\cite{FoSt18},
investigated the distribution of the discrete spectrum of
the Laplacian $\square_{M}$ of some pseudo-Riemannian manifolds, i.e., when
$M$ is the flat pseudo-Riemannian manifold $\mathbb{R}^{p,q}/\mathbb{Z}^{p+q}$
and is the Lorentzian manifold $S^1\times S^{q}$, respectively.

Let us recall some basic notions.
A \textit{discontinuous group} for a homogeneous manifold \mbox{$X\!=\!G/H$}
is a discrete subgroup $\Gamma$ of $G$ acting properly discontinuously
and freely on $X$ (Kobayashi \cite[Definition~1.3]{Ko01}). In~this case, the quotient space $X_{\Gamma}:=\Gamma\backslash X$ carries a
$C^{\infty}$-manifold structure such that
the quotient map $p_{\Gamma}\colon X\rightarrow X_{\Gamma}$ is a covering of
$C^{\infty}$ class, hence $X_{\Gamma}$ has a $(G,X)$-structure induced by $p_{\Gamma}$.
If we drop the assumption of freeness, $X_{\Gamma}$ is not always a manifold but
carries a nice structure called an orbifold or $V$-manifold.
Proper discontinuity is a more
serious assumption which assures $X_{\Gamma}$ to be Hausdorff in the quotient topology. We~remark that the action of a discrete subgroup $\Gamma$ on $X$
may fail to be properly discontinuous
when $H$ is noncompact. In~order to overcome this difficulty,
Kobayashi~\cite{Kob96} and Benoist~\cite{Ben96} established the properness criterion
for reductive $G$ generalizing the original criterion by Kobayashi~\cite{Kob89}.
Whereas discontinuous groups for the de Sitter space
$\mathrm{dS}^{n}:=\mathrm{SO}_{0}(n,1)/\mathrm{SO}_{0}(n-1,1)$
are always finite groups
(the Calabi--Markus phenomenon, see~\cite{CaMa62, Kob89}),
there are a rich family of discontinuous groups for
the anti-de Sitter space, see, e.g.,~\cite{Go85, Kob98, Sa00}. We~treat, in this article, the three-dimensional anti-de Sitter space
$\mathrm{AdS}^{3}:=\mathrm{SO}_{0}(2,2)/(\{\pm1\}\times\mathrm{SO}_{0}(2,1))$.

For $m\in \mathbb{N}$, we set
\begin{gather*}%\label{lambda}
\lambda_{m} := 4m(m-1).
\end{gather*}


% Original source lines 242--472
\section{Preliminaries about the anti-de Sitter space}
\label{Preliminary-AdS}

In this section, we collect some preliminary results about $\mathrm{AdS}^{3}$. We~refer to \cite[Section~9]{KaKob16} where they illustrate their general theory
for reductive symmetric spaces $X=G/H$ in details in the special setting where
$X=\mathrm{AdS}^{3}$. See also~\cite{Kan19}.

Let $Q$ be a quadratic form on $\mathbb{R}^{4}$ defined by
$Q(x) = x_1^2 + x_2^2-x_3^2-x_4^2$ for $x= (x_1,x_2,x_3,x_4)$ and
we set
\begin{gather*}
\mathbb{H}^{2,1}:=\big\{x=(x_1,x_2,x_3,x_4) \in \mathbb{R}^4 \mid Q(x) = 1\big\}
\cong \mathrm{SO}_{0}(2,2)/\mathrm{SO}_{0}(2,1).
\end{gather*}
The tangent space $T_{x}(\mathbb{H}^{2,1})$ at $x\in\mathbb{H}^{2,1}$
is isomorphic to the orthogonal complement $(\mathbb{R} x)^{\bot}$ with respect to $Q$.
Then $-Q|_{(\mathbb{R} x)^{\bot}}$ is
a quadratic form of signature $(2,1)$ on
$T_{x}(\mathbb{H}^{2,1}) \cong (\mathbb{R} x)^{\bot}$ and
thus $-Q$ induces a Lorentzian structure on $\mathbb{H}^{2,1}$
with constant sectional curvature $-1$.
The~$3$-di\-men\-sional anti-de Sitter space
\begin{gather*}
\mathrm{AdS}^{3}:=\mathbb{H}^{2,1}/\{\pm 1\}
\cong \mathrm{SO}_{0}(2,2)/(\{\pm1\}\times\mathrm{SO}_{0}(2,1)),
\end{gather*}
inherits a Lorentzian structure through
the double covering $\pi\colon\mathbb{H}^{2,1}\rightarrow\mathrm{AdS}^{3}$.

\subsection{Some coordinates and ``pseudo-balls''}
%\label{coord}
In this subsection, we work with coordinates on $\mathbb{H}^{2,1}$ and
consider ``pseudo-balls'' in $\mathrm{AdS}^{3}$. We~identify
$\mathbb{H}^{2,1}$ with $\mathrm{SL}(2,\mathbb{R})$ using the isomorphism
\begin{align}
\begin{array}{@{}ccc}
\mathbb{H}^{2,1} & \xrightarrow{\cong} & \mathrm{SL}(2,\mathbb{R}), \\
x=(x_1,x_2,x_3,x_4) & \longmapsto & \begin{pmatrix}
x_1 + x_4 & -x_2 + x_3 \\
x_2 + x_3 & x_1 -x_4
\end{pmatrix}\!.\label{isom}
\end{array}
\end{align}
For $t\geq0$ and $\theta\in\mathbb{R}$, we use the notations
\begin{gather}
\label{k-a}
k(\theta) = \begin{pmatrix}
\cos\theta & -\sin\theta \\
\sin\theta & \cos\theta
\end{pmatrix}\!,\qquad
a(t) = \begin{pmatrix}
{\rm e}^t & 0 \\
0 & {\rm e}^{-t}
\end{pmatrix}\!.
\end{gather}
We embed $\mathbb{H}^{2,1}$ into $\mathbb{C}^{2}$ by
\begin{gather}
\label{embed}
x\mapsto (z_1,z_2) = \big(x_1 + \sqrt{-1}x_2, x_3 + \sqrt{-1}x_4\big).
\end{gather}
We note that $z_{1}\neq0$ if $x\in\mathbb{H}^{2,1}$.
Via the identification (\ref{isom}), we have
\begin{gather}
\label{polar}
(z_1,z_2) = \big((\cosh t){\rm e}^{\sqrt{-1}(\theta_1+\theta_2)},
(\sinh t){\rm e}^{\sqrt{-1}(\theta_1-\theta_2)}\big),
\end{gather}
if $x = k(\theta_1)a(t)k(\theta_2) \in \mathrm{SL}(2,\mathbb{R})$
(a ``polar coordinate''). In~particular, we have
\begin{gather*}
\cosh 2t = x_1^2 + x_2^2 + x_3^2 + x_4^2.
\end{gather*}
Next, we consider pseudo-balls on $\mathrm{AdS}^3$, as a special case of
Kassel--Kobayashi~\cite{KaKob16} for
reductive symmetric spaces.
\begin{Definition}
\label{def-pball}
For $x = (x_1,x_2,x_3,x_4)\in\mathbb{H}^{2,1}$, $\|x\|\in\mathbb{R}_{\geq0}$
is defined by
\begin{gather*}
\cosh\|x\|:=x_1^2 + x_2^2 + x_3^2 + x_4^2\qquad (= \cosh(2t)).
\end{gather*}
This function is invariant under $x\mapsto -x$, hence defines
a function on $\mathrm{AdS}^{3}$,
to be also denoted by $\|\cdot\|$ (a ``pseudo-distance'' from the origin).
The compact set
\begin{gather*}
B(R) := \big\{y\in\mathrm{AdS}^{3} \mid \|y\| \leq R\big\}
\end{gather*}
is called a pseudo-ball of radius $R$.
\end{Definition}

\subsection[Square-integrable eigenfunctions of the Laplacian on the anti-de Sitter space] {Square-integrable eigenfunctions of the Laplacian \\on the anti-de Sitter space}\label{preliminary-spherical}

In this subsection, we consider square-integrable eigenfunctions of
$\square_{\mathrm{AdS}^{3}}$ with eigenvalues $\lambda_{m}=4m(m-1)$. We~recall from \cite[Section~9]{KaKob16} the following decomposition of the open subset $\{Q > 0\}$ of
the flat pseudo-Riemannian manifold $\mathbb{R}^{2,2}=\big(\mathbb{R}^{4},Q({\rm d}x)\big)$:
\begin{align*}
\begin{array}{@{}ccc}
\{Q > 0\} & \xrightarrow{\cong} & \mathbb{R}_{>0} \times \mathbb{H}^{2,1}, \\[1ex]
x & \longmapsto & \big(\sqrt{Q(x)}, x/\sqrt{Q(x)}\big).
\end{array}
\end{align*}
Let $r=\sqrt{Q(x)}$. Then one has, see \cite[p.~215]{KaKob16},
\begin{align}
\label{decomposition-of-laplacian}
-r^2\square_{\mathbb{R}^{2,2}} = -\left(r\frac{\partial}{\partial r}\right)^2 - 2r\frac{\partial}{\partial r} + \square_{\mathbb{H}^{2,1}}.
\end{align}
Let $m$ be a positive integer and $k$ be a non-negative integer. In~the coordinates (\ref{embed}),
the homogeneous function $z_{1}^{-(k+2m)}z_{2}^{k}$ of degree $-2m$
 is harmonic with respect to $\square_{\mathbb{R}^{2,2}}$,
hence its restriction to the submanifold $\mathbb{H}^{2,1}$ is
an eigenfuction of $\square_{\mathbb{H}^{2,1}}$ with eigenvalue $\lambda_{m}=4m(m-1)$
by~the formula (\ref{decomposition-of-laplacian}).
Moreover, it is square-integrable with respect to the measure
$\sinh(2t){\rm d}\theta_{1}{\rm d}t{\rm d}\theta_{2}$ in~the polar coordinate (\ref{polar})
induced from the Lorentzian metric on $\mathbb{H}^{2,1}$,
as in the $k=0$ case \cite[Section~9]{KaKob16}.
This $L^{2}$-eigenfunction is invariant under $(z_{1},z_{2})\mapsto(-z_{1},-z_{2})$,
hence defines a~real analytic $L^{2}$-eigenfunction on $\mathrm{AdS}^{3}$
with eigenvalue $\lambda_{m}$, to be denoted by $\psi_{m,k}$.
The discrete spectrum
$\mathrm{Spec}_{d}(\square_{\mathrm{AdS}^{3}})$ coincides with
$\{\lambda_{m} \mid m \in \mathbb{N}\}$
and $L^{2}_{\lambda_{m}}\big(\mathrm{AdS}^{3}\big)$ is generated by
$\psi_{m,0}$ and its complex conjugate $\overline{\psi_{m,0}}$
as a representation of $\mathrm{SO}_{0}(2,2)$ (see \cite[Claim~9.12]{KaKob16}).
By~(\ref{polar}), we~have
\begin{align}
\label{polarFJ}
\psi_{m,k}(\pi(x))
= {\rm e}^{-2\sqrt{-1}(m\theta_1 + (m+k)\theta_2)}\tanh^{k}t\cosh^{-2m}t
\end{align}
for $x=k(\theta_1)\,a(t)\,k(\theta_2)\in\mathbb{H}^{2,1}$. We~refer to $\psi_{m,k}$ as a spherical function of type $(-m,m+k)$
in~accordance with the action of
$\mathrm{SO}(2)\times\mathrm{SO}(2)$.

\subsection{Convergence of generalized Poincar\'e series}\label{preliminary-Poincare}
In this subsection, we explain the fact about the discrete spectrum
of locally symmetric spaces by Kassel--Kobayashi~\cite{KaKob16}
in our $\mathrm{AdS}^{3}$ setting. We~use the following notation.
\begin{Notation}\quad
\begin{itemize}\itemsep=0pt
\item
Let $\grave{\ }G=\mathrm{PSL}(2,\mathbb{R})=
\mathrm{SL}(2,\mathbb{R})/\{\pm1\}$ and
$G=\grave{\ }G\times\grave{\ }G$.
\item
Let $\grave{\ }K=\mathrm{PSO}(2)=\mathrm{SO}(2)/\{\pm1\}$ and
$K = \grave{\ }K\times\grave{\ }K$.
\item
Let $E$ and $\grave{\ }E$ be respectively the identity elements of $G$ and $\grave{\ }G$.
\end{itemize}
\end{Notation}

\begin{Remark}
The double covering $\mathrm{SO}_{0}(2,2)\rightarrow G$ induces
an isomorphism $\mathrm{AdS}^{3}\cong G/\mathrm{diag} \grave{\ }G$ $(\cong \grave{\ }G)$.
From now on, we consider only discontinuous groups $\Gamma$ for $\mathrm{AdS}^{3}$
which are discrete subgroups of $G$.
This is enough for our purpose.
\end{Remark}

In order to study
$\operatorname{Spec}_{d}\big(\square_{\Gamma\backslash\mathrm{AdS}^3}\big)$,
Kassel--Kobayashi~\cite{KaKob16} considered the convergence and non-vanishing of
generalized Poincar\'e series
\begin{gather}
\label{Poincare-series}
\varphi^{\Gamma}(\Gamma x):=\sum_{\gamma \in \Gamma} \varphi\big(\gamma^{-1}x\big)
\end{gather}
for $K$-finite square-integrable eigenfunctions $\varphi$ of $\square_{\mathrm{AdS}^{3}}$.
For this, they used an analytic estimate of $\varphi$ and
a geometric estimate of the number of $\Gamma$-orbits
\begin{gather}
\label{orbit-count}
N_{\Gamma}(x,R) := \#\{\gamma \in \Gamma \mid \gamma x \in B(R)\}
\end{gather}
in the pseudo-ball $B(R)$ for $R>0$.
Since the $\Gamma$-action is properly discontinuous and $B(R)$ is compact,
we have $N_{\Gamma}(x,R)<\infty$.

The convergence of generalized Poincar\'e series is proved by~\cite{KaKob16} as follows.
For $g \in G$ and a~function $f$ on $\mathrm{AdS}^{3}$,
$\ell_{g}^{*}f$ is defined by $\ell_{g}^{*}f(x)=f\big(g^{-1}x\big)$.
\begin{Fact}[Kassel--Kobayashi~\cite{KaKob16}]\label{Poincare}
Let $\Gamma \subset G$ be a discontinuous group for $\mathrm{AdS}^{3}$
satisfying the exponential growth condition
\begin{gather}\label{C}
\exists A,a > 0,\qquad
\forall x\in \mathrm{AdS}^{3},\qquad
\forall R>0, \qquad
N_{\Gamma}(x,R) < A{\rm e}^{aR}.
\end{gather}
Then, for any $K$-finite eigenfunction $\varphi$ of $\square_{\mathrm{AdS}^{3}}$ with eigenvalue
$\lambda_{m}$
and any $g\in G$,
if $m > a$, then
$(\ell_{g}^{*}\varphi)^{\Gamma}$ $($see $\eqref{Poincare-series})$ is continuous and square-integrable on
$\Gamma\backslash\mathrm{AdS}^{3}$ and
an eigenfunction of~$\square_{\Gamma\backslash\mathrm{AdS}^{3}}$ with eigenvalue
$\lambda_{m}$.
\end{Fact}

\begin{Remark}\label{remPoincare}\quad
\begin{enumerate}\itemsep=0pt
\item[(1)]
Fact~\ref{Poincare} does not assert the non-vanishing of the series
$(\ell_{g}^{*}\varphi)^{\Gamma}$ which is more involved.
Kassel--Kobayashi~\cite{KaKob16} proved that
there exists $g\in G$ such that
$(\ell^{*}_{g}\psi_{m,0})^{\Gamma}\neq 0$
for sufficiently large $m \in \mathbb{N}$.
\item[(2)]
%\label{remcount}
By \cite[Lemma~4.6.4]{KaKob16}, if a discontinuous group $\Gamma$ is sharp
in the sense of \cite[Definition~4.2]{KaKob16},
then $\Gamma$ satisfies the exponential growth condition (\ref{C}).
Moreover, Kassel~\cite{Ka09} and Gu\'eriataud--Kassel~\cite{GuKa17}
proved that finitely generated discontinuous groups for $\mathrm{AdS}^{3}$
are always sharp (see Fact~\ref{proper} below).
\item[(3)]
There exist discontinuous groups
which do not satisfy the exponential growth condition~(\ref{C}).
Indeed, for any increasing function $f\colon\mathbb{R}\rightarrow \mathbb{R}_{>0}$
and any $x\in\mathrm{AdS}^{3}$,
we~con\-st\-ructed
a discontinuous group $\Gamma_{f,x}$ for $\mathrm{AdS}^{3}$
satisfying $N_{\Gamma_{f,x}}(x,R) > f(R)$ for sufficiently large $R>0$ in~\cite{Kan19}.
\end{enumerate}
\end{Remark}

\setcounter{Theorem}{5}\setcounter{subsection}{3}\setcounter{equation}{9}
% Original source lines 498--845
\subsection{``Injectivity radii'' of anti-de Sitter 3-manifolds}\label{inj-radii}
Let $\Gamma$ be a discontinuous group for $\mathrm{AdS}^{3}$. In~this subsection, we give a uniform estimate of the pseudo-distance
between the origin and the second closest point of each $\Gamma$-orbit.

We recall that $\Gamma(\subset \grave{\ }G\times\grave{\ }G)$
acts isometrically on $\mathrm{AdS}^{3}(\cong\grave{\ }G)$ by
$(\gamma_{1},\gamma_{2})x=\gamma_{1} x\gamma_{2}^{-1}$
for~$(\gamma_{1},\gamma_{2})\in\Gamma$ and $x\in\grave{\ }G$. We~set
\begin{align}
\label{varepsilon}
\varepsilon_{\Gamma} := \inf_{(\gamma_1,\gamma_2)\in
\Gamma\setminus\{E\}}
\frac{1}{3}\big|\|\gamma_1\|-\|\gamma_2\|\big|.
\end{align}
By the inequality (see, e.g., \cite[Lemma~5.5]{Kan19})
\begin{gather*}
\|(g_1,g_2)x\|\geq \big | \|g_1\|-\|g_2\| \big | - \|x\|\qquad
\text{for}\quad (g_1,g_2) \in G \quad\text{and}\quad x \in \mathrm{AdS}^{3},
\end{gather*}
we get:
\begin{Lemma}\label{inj-rad}
If $\varepsilon_{\Gamma}>0$, then
$\gamma B(\varepsilon_{\Gamma}) \cap B(\varepsilon_{\Gamma})
= \varnothing$
for all $\gamma\in \Gamma\setminus\{E\}$.
\end{Lemma}

\begin{Proposition}\label{strong}
Let $\Gamma$ be a discrete subgroup of $G$ acting properly discontinuously
on $\mathrm{AdS}^{3}$.
Then there exists $g \in G$
satisfying $\varepsilon_{g^{-1}\Gamma g} > 0$.
\end{Proposition}
\begin{Remark}
One sees in the proof below that
the set of such $g$ is dense in $G$.
\end{Remark}
Proposition~\ref{strong} follows obviously
from the proper discontinuity of the $\Gamma$-action and the following lemma:
\begin{Lemma}\label{weak}
For any countable subset $\Gamma$ of
$G$,
there exists
$g\in G$
such that $\|\gamma_1\|\neq\|\gamma_2\|$
for all $\gamma=(\gamma_1,\gamma_2)\in
g^{-1}\Gamma g \setminus \{E\}$.
\end{Lemma}

\begin{proof}[Proof of Lemma~\ref{weak}]
For $\gamma \in \Gamma$,
the map $f_{\gamma}\colon G \rightarrow G$ defined by $g\mapsto g^{-1}\gamma g$
is real analytic.
For the analytic subset
$F = \{(g_1,g_2) \in G
\mid \|g_1\| = \|g_2\| \}$ of $G$, we claim that
the set $f_{\gamma}^{-1}(F)$ is a proper subset of $G$ if $\gamma\neq E$.
For this, we may assume $\gamma_1 \neq \grave{\ }E$ without loss of generality.
Then there exists
$g_1 \in \grave{\ }G$ satisfying $\|g_1^{-1}\gamma_1g_1\| \neq \|
\gamma_1\|$ as one can find $g_{1}$ depending on the three cases where
$\gamma_{1}$ is hyperbolic, parabolic, or elliptic.
Hence $(g_1,\grave{\ }E)\notin f_{\gamma}^{-1}(F)$ if $\|\gamma_{1}\|=\|\gamma_{2}\|$,
and $E\notin f_{\gamma}^{-1}(F)$ if not. Thus $f_{\gamma}^{-1}(F)$ is a proper subset of $G$.

Therefore the analytic set
$f_{\gamma}^{-1}(F)$ has no interior point, and thus so does the countable union
$\bigcup_{\gamma \in \Gamma \setminus \{E\}}
 f_{\gamma}^{-1}(F)$ by the Baire category theorem
(see, e.g., \cite[Theorem~2.2]{functional-analysis}).
Hence there exists an element
$g$ of $G\setminus\bigcup_{\gamma \in \Gamma \setminus \{E\}}
 f_{\gamma}^{-1}(F)$ and we have
$\|\gamma_1\|\neq\|\gamma_2\|$
for all $\gamma=(\gamma_1,\gamma_2)\in
g^{-1}\Gamma g \setminus \{E\}$.
\end{proof}

\section{Proof of Theorem~\ref{mthm1}}\label{proof-linear-indep}
In this section, we prove Theorem~\ref{mthm1}.

More generally, without finitely generated assumption of $\Gamma$,
we study linear independence of~the generalized Poincar\'e series
of the spherical functions $\psi_{m,k}$ of type $(-m,m+k)$
defined in~Section~\ref{preliminary-spherical}.
By choosing $k=3^{j}$ ($j=0,1,2,\ldots$),
we prove:
\begin{Theorem}
\label{unbound}
If $\Gamma$ is a discontinuous group for $\mathrm{AdS}^{3}$ satisfying
the exponential growth condition~\eqref{C}, then
\begin{gather*}
\lim_{m\to\infty}\mathcal{N}_{\Gamma\backslash\mathrm{AdS}^{3}}(\lambda_{m})=\infty.
\end{gather*}
\end{Theorem}
Theorem~\ref{mthm1} is a direct consequence of
Theorem~\ref{unbound} by Remark~\ref{remPoincare}(2).

\begin{Proposition}
\label{mainthm}
Let $\Gamma$ be a discrete subgroup of $G$ acting properly discontinuously
on $\mathrm{AdS}^{3}$
and satisfying the exponential growth condition \eqref{C}.
If $\varepsilon_{\Gamma}>0$,
then there exists a real number $m_{\Gamma}(k)$
$($given explicitly by $\eqref{Poincare_estimate})$
for $k\in \mathbb{N}$
such that
$\{(\operatorname{Re}(\psi_{m,3^{j}}))^{\Gamma}\}_{j=0}^{k-1}
\subset L^2_{\lambda_{m}}(\Gamma\backslash\mathrm{AdS}^{3})$
are linearly independent
for all integers $m > m_{\Gamma}(k)$.
\end{Proposition}

Postponing the proof of Proposition~\ref{mainthm} until the end of this section,
we prove Theorem~\ref{unbound}.

\begin{proof}[Proof of Theorem~\ref{unbound}]
We have an obvious equality of the multiplicity of $L^2$-eigenvalues,
$\mathcal{N}_{\Gamma\backslash\mathrm{AdS}^{3}} =
\mathcal{N}_{(g^{-1}\Gamma g)\backslash\mathrm{AdS}^{3}}$
for any $g\in G$
through the natural isomorphism
$\Gamma\backslash\mathrm{AdS}^{3} \cong
\big(g^{-1}\Gamma g\big)$ $\backslash\mathrm{AdS}^{3}$ as Lorentzian manifolds.
By replacing $\Gamma$ with $g^{-1}\Gamma g$ if necessary,
we may and do assume $\varepsilon_{\Gamma}>0$ by Proposition~\ref{strong}.
Then Proposition~\ref{mainthm} implies that
$L^2_{\lambda_{m}}\big(\Gamma\backslash \mathrm{AdS}^{3}\big)$ contains at least
$k$ linearly independent elements
if $m>m_{\Gamma}(k)$ for any fixed $k\in\mathbb{N}$,
which means
$\dim_{\mathbb{C}} L^2_{\lambda_{m}}(\Gamma\backslash\mathrm{AdS}^{3})
\geq k.$
Hence Theorem~\ref{unbound} follows.
\end{proof}

Kassel--Kobayashi~\cite{KaKob16} proved the non-vanishing of
the generalized Poincar\'e series $(\psi_{m,0})^{\Gamma}$
for sufficiently large $m\in\mathbb{N}$
by showing that the first term in the generalized Poincar\'e series
is larger at the origin than the sum of the remaining terms.
For this, they utilized the fact that $\psi_{m,0}(\grave{\ }E)=1$.
Our strategy for the proof of Proposition~\ref{mainthm} is along the same line,
however, there are some technical difficulties
since $\psi_{m,k}$ for $k\geq1$ vanishes at the origin. We~then make use of an observation that
$\psi_{m,k}$ decays more slowly at the origin than at infinity,
to be precise, by the following formula, see (\ref{polarFJ}):
\begin{gather*}
|\psi_{m,k}(x)|=\cosh^{-2m}(\|x\|/2) \tanh^{k}(\|x\|/2).
\end{gather*}
Actually, we use an analytic lemma (Lemma~\ref{koukou})
to prove that the first term in the generalized Poincar\'e series $(\psi_{m,k})^{\Gamma}$
is larger at points sufficiently close to the origin
than the sum of the remaining terms if $m\gg0$.
Moreover,
we use a combinatorial lemma (Lemma~\ref{sekoi})
to find points at which
leading terms of
$(\operatorname{Re}(\psi_{m,k}))^{\Gamma}$
do not cancel each other for any linear combination.

For $C,a,\varepsilon > 0$ and $s \in \mathbb{N}$, we set
\begin{align*}%\label{inequality}
m(C,a,\varepsilon,s) := \frac{(\log2)s + 2a\varepsilon + \log\big(1+2^{s}C{\rm e}^{6a\varepsilon}\big)}{\log\cosh\varepsilon}
\end{align*}
and
\begin{align*}
\tilde{m}(C,a,\delta,s) :=
\inf_{0<\varepsilon<\delta}m(C,a,\varepsilon,s).
\end{align*}
Note that $\tilde{m}(C,a,\delta,s)=O\big(\delta^{-2}\big)$ as $\delta\to 0$
and $=O(1)$ as $\delta\to \infty$.
\begin{Lemma}\label{koukou}
For any integer $m > m(C,a,\varepsilon,s)$
and any one-variable polynomial $f$ of degree $\leq s$ with non-negative coefficients,
\begin{gather*}
C\sum_{n=1}^{\infty}{\rm e}^{4a(n+1)\varepsilon}(\cosh 2n\varepsilon)^{-m}f(\tanh
2(n+1)\varepsilon) < (\cosh \varepsilon)^{-m}f(\tanh \varepsilon).
\end{gather*}
\end{Lemma}
\begin{proof}
We may assume that $f(x) = x^{j}$ for $j=0,1,\ldots,s$.
Since
\begin{gather*}
1\leq \frac{\tanh nx}{\tanh x} \leq n,\
(\cosh x)^{n} \leq \cosh nx
\end{gather*}
for $x\in\mathbb{R}$, we have
\begin{align*}
\text{(LHS)/(RHS)} &=
C\sum_{n=1}^{\infty}{\rm e}^{4a(n+1)\varepsilon}\bigg(\frac{\cosh 2n
\varepsilon}{\cosh \varepsilon}\bigg)^{-m}\bigg(\frac{\tanh
2(n+1)\varepsilon}{\tanh \varepsilon}\bigg)^j \nonumber
\\
&\leq C{\rm e}^{6a\varepsilon}\sum_{n=1}^{\infty}
\big({\rm e}^{2a\varepsilon}(\cosh \varepsilon)^{-m}\big)^{2n-1}(2(n+1))^s.
\end{align*}
We set $d:={\rm e}^{2a\varepsilon}(\cosh \varepsilon)^{-m}$.
Then $d<1$ by $m>m(C,a,\varepsilon,s)$. Since $n+1\leq 2^{n}$ for all $n\in\mathbb{N}$,
we have
\begin{gather*}
\text{(LHS)/(RHS)}\leq
2^{s}C{\rm e}^{6a\varepsilon}\sum_{n=1}^{\infty} (2^{s}d)^n
= 2^{s}C{\rm e}^{6a\varepsilon}\frac{2^{s}d}{1-2^{s}d}.
\end{gather*}
Again by $m > m(C,a,\varepsilon,s)$,
we have $2^{s}d<\big(1+2^{s}C{\rm e}^{6a\varepsilon}\big)^{-1}$.
Therefore we obtain
\begin{gather*}
\text{(LHS)/(RHS)}<1.\tag*{\qed}
\end{gather*}
\renewcommand{\qed}{}
\end{proof}
Let $\chi\colon\{\pm 1\} \rightarrow \{0,1\}$ be
the map defined by $\chi(1) = 0$ and $\chi(-1) = 1$.
For $a = (a_j)_{j=0}^{k-1} \in \{\pm1\}^{k}$ and an odd integer $N\geq 3$,
we set
\begin{align*}%\label{theta}
\theta_{a,N} := \pi\sum_{i=0}^{k-1} (\chi(a_i) - \chi(a_{i-1}))N^{-i}.
\end{align*}
Here we use the convention $a_{-1}=1$.
\begin{Lemma}\label{sekoi}
For any $a=(a_{0},\ldots,a_{k-1})\in \{\pm1\}^{k}$ and any odd integer $N$,
we have
\begin{gather*}
a_j\cos\big(N^{j}\theta_{a,N}\big) > 0 \qquad \text{for}\quad j=0,1,\ldots,k-1.
\end{gather*}
\end{Lemma}

\begin{proof}
Since
$N^{k-1}\theta_{a,N} \equiv \pi \chi(a_{k-1}) \ (\mathrm{mod}\ 2\pi)$,
we have $\cos \big(N^{k-1} \theta_{a,N}\big) = a_{k-1}$.
It is easy to check that
$|N^{j} \theta_{a,N} - N^{j} \theta_{(a_0,\cdots,a_j),N}| < \pi/2$ for $j=0,1,\ldots,k-1$,
hence the signature of~$\cos (N^{j} \theta_{a,N})$ is equal to that of
$\cos (N^{j} \theta_{(a_0,\cdots,a_j),N}) = a_j$.
\end{proof}

\begin{Remark}
We have used the geometric progression $(N^{j})_{j=0}^{k-1}$ in Lemma~\ref{sekoi}.
On the other hand, an analogous statement does not hold if
we use arithmetic progressions. For example,
there does not exist $\theta \in \mathbb{R}$
satisfying $a_j\cos m_j\theta > 0$ for all $j=0,1,2,3,4$
if we choose $(a_j)_{j=0}^{4}=(1,1,1,-1,1)$ and
an arithmetic progression $(m_j)_{j=0}^{4}$.
\end{Remark}

For a discontinuous group $\Gamma$ and $k\in \mathbb{N}$,
one can take $m_{\Gamma}(k)$ in Proposition~\ref{Poincare_estimate} by
\begin{gather}
\label{Poincare_estimate}
m_{\Gamma}(k)=\inf_{(A,a)\in \mathrm{C_{exp}}(\Gamma)}
\max\big\{\tilde{m}\big(3^{k-1}A,a,\varepsilon_{\Gamma}/4,3^{k-1}\big)/2,a\big\},
\end{gather}
where $\mathrm{C_{exp}}(\Gamma):=\big\{(A,a)\in \mathbb{R}^{2} \mid
\forall x\in \mathrm{AdS}^{3},\forall R>0, N_{\Gamma}(x,R) < A{\rm e}^{aR}\big\}$.
Here,
we adopt the convention that $\inf_{\varnothing} f =\infty$ for a real-valued function $f$. In~particular,
$m_{\Gamma}(k)=\infty$ when $\mathrm{C_{exp}}(\Gamma)=\varnothing$
or $\varepsilon_{\Gamma}=0$.
\begin{proof}[Proof of Proposition~\ref{mainthm}]
By the exponential growth condition (\ref{C}), $\mathrm{C_{exp}}(\Gamma)\neq\varnothing$
and thus $m_{\Gamma}(k)<\infty$. We~take an integer $m >m_{\Gamma}(k)$.
Then there exist $\varepsilon$ with $0<\varepsilon<\varepsilon_{\Gamma}/4$
and $(A,a)\in\mathrm{C_{exp}}(\Gamma)$
satisfying the inequality
$m > \max\big\{m\big(3^{k-1}A,a,\varepsilon,3^{k-1}\big)/2,a\big\}$.

To see $\mathbb{C}$-linear independence of the real-valued functions
$\big\{(\operatorname{Re}(\psi_{m,3^{j}}))^{\Gamma}\big\}_{j=0}^{k-1}$,
it is enough to prove the non-vanishing of
the real part $\operatorname{Re}\big(\psi_{m,b}^{\Gamma}\big)
=(\operatorname{Re}(\psi_{m,b}))^{\Gamma}$ of
the generalized Poincar\'e series of a linear combination
\begin{gather*}
\psi_{m,b} := \sum_{j = 0}^{k-1} b_j\psi_{m,3^{j}}
\end{gather*}
for any
$b=(b_{0},b_{1},\ldots,b_{k-1}) \in \mathbb{R}^{k}\setminus\{0\}$.
By Lemma~\ref{inj-rad}, for $x\in B(4\varepsilon)$, we have
\begin{gather}
\label{separate}
\psi_{m,b}^{\Gamma}(\Gamma x)
= \psi_{m,b}(x) + \sum_{\substack{\gamma \in \Gamma \\
\|\gamma^{-1} x\| > 4\varepsilon}} \psi_{m,b}\big(\gamma^{-1} x\big).
\end{gather}
By (\ref{polarFJ}), for any $y \in \mathrm{AdS}^{3}$, we get
\begin{gather*}
|\psi_{m,b}(y)| \leq \bigg( \!\cosh \frac{\| y \|}{2} \bigg)^{-2m}
\sum_{j=0}^{k-1} |b_{j}|\bigg(\! \tanh \frac{\| y \|}{2} \bigg)^{3^{j}}.
\end{gather*}
We define
$a=(a_{j})_{j=0}^{k-1}$ by $a_{j}=1$ for $b_{j}\geq 0$ and
$a_{j}=-1$ for $b_{j}<0$,
and set
\begin{gather*}
%\label{pol-b}
f_{b}(u) := \sum_{j=0}^{k-1} b_{j}\cos\big(3^{j}\theta_{a,3}\big)u^{3^{j}}.
\end{gather*}
We note that all the coefficients of $f_{b}$ are non-negative by Lemma~\ref{sekoi}.
Moreover, we get
$\big| \cos \big(3^{j}\theta_{a,3}\big) \big|^{-1}
\leq 3^{k-1}$ for all $j=0,1,\ldots,k-1$
by using the inequality $\sin(\pi x/2)\geq x$ for~$0\leq x\leq1$.
Thus
\begin{gather*}
|\psi_{m,b}(y)| \leq 3^{k-1}\bigg(\!\cosh \frac{\| y \|}{2} \bigg)^{-2m}
f_{b}\bigg(\! \tanh \frac{\| y \|}{2} \bigg)
\end{gather*}
and, for any $x\in B(4\varepsilon)$, we have
\begin{align}
\bigg|\!\!\!\sum_{\substack{\gamma \in \Gamma \\ \|\gamma^{-1} x\| > 4\varepsilon}} \operatorname{Re}\big(\psi_{m,b}\big(\gamma^{-1} x\big)\big)\bigg|
&\leq \sum_{n = 1}^{\infty} \sum_{\substack{\gamma \in \Gamma \\
4\varepsilon n < \|\gamma^{-1} x\| \leq 4\varepsilon (n + 1) }}
|\psi_{m,b}(\gamma^{-1} x)| \nonumber \\
&\leq 3^{k-1}\sum_{n = 1}^{\infty} N_{\Gamma}(x, 4\varepsilon (n + 1))\left(\cosh
2\varepsilon n \right)^{-2m} f_{b}\left(\tanh 2\varepsilon (n+1)\right)
\nonumber \\
&\leq 3^{k-1}A\sum_{n = 1}^{\infty} {\rm e}^{4a\varepsilon (n+1)}\left(\cosh 2\varepsilon n
\right)^{-2m} f_{b}\left(\tanh 2\varepsilon (n+1)\right)
\nonumber \\
&< \left(\cosh \varepsilon\right)^{-2m} f_{b}
\left(\tanh \varepsilon\right).
\label{small}
\end{align}
The third and forth inequalities respectively follow from
the exponential growth condition (\ref{C}) and Lemma~\ref{koukou}.
On the other hand, we set
\begin{gather*}
%\label{non-zero-point}
x_{a,\varepsilon}:=
k\bigg(\frac{\theta_{a,3}}{2}\bigg)a(\varepsilon)k\bigg(\frac{\theta_{a,3}}{2}\bigg)^{-1} \in B(4\varepsilon).
\end{gather*}
Then it follows from (\ref{polarFJ}) that
\begin{align}
\label{large}
\operatorname{Re} \psi_{m,b}(x_{a,\varepsilon})
=\left(\cosh \varepsilon\right)^{-2m} f_{b}
\left(\tanh \varepsilon\right).
\end{align}
By (\ref{separate}), (\ref{small}), and (\ref{large}),
we obtain $(\operatorname{Re}(\psi_{m,b}))^{\Gamma}(\Gamma x_{a,\varepsilon})\neq 0$.
Hence we complete the proof
by the continuity of $\psi_{m,b}^{\Gamma}$ (Fact~\ref{Poincare}).
\end{proof}
\begin{thebibliography}{99}
\bibitem[1]{Ben96}
Benoist Y., Actions propres sur les espaces homog\`enes r\'eductifs,
 \href{https://doi.org/10.2307/2118594}{\textit{Ann. of Math.}} \textbf{144} (1996), 315--347.

\bibitem[3]{CaMa62}
Calabi E., Markus L., Relativistic space forms, \href{https://doi.org/10.2307/1970419}{\textit{Ann. of Math.}}
 \textbf{75} (1962), 63--76.

\bibitem[4]{FoSt18}
Fox J., Strichartz R.S., Unexpected spectral asymptotics for wave equations on
 certain compact spacetimes, \href{https://doi.org/10.1007/s11854-018-0059-2}{\textit{J.~Anal. Math.}} \textbf{136} (2018),
 209--251, \href{https://arxiv.org/abs/1407.2517}{arXiv:1407.2517}.

\bibitem[5]{Go85}
Goldman W.M., Nonstandard {L}orentz space forms, \href{https://doi.org/10.4310/jdg/1214439567}{\textit{J.~Differential Geom.}}
 \textbf{21} (1985), 301--308.

\bibitem[6]{GuKa17}
Gu\'eritaud F., Kassel F., Maximally stretched laminations on geometrically
 finite hyperbolic manifolds, \href{https://doi.org/10.2140/gt.2017.21.693}{\textit{Geom. Topol.}} \textbf{21} (2017),
 693--840, \href{https://arxiv.org/abs/1307.0250}{arXiv:1307.0250}.

\bibitem[7]{Kan19}
Kannaka K., Counting orbits of certain infinitely generated non-sharp
 discontinuous groups for the anti-de {S}itter space, \href{https://arxiv.org/abs/1907.09303}{arXiv:1907.09303}.

\bibitem[8]{Ka09}
Kassel F., Quotients compacts d'espaces homog\`enes r\'eels ou $p$-adiques,
 Ph.D.~Thesis, {U}niversit\'e Paris-Sud, 2009.

\bibitem[10]{KaKob16}
Kassel F., Kobayashi T., Poincar\'e series for non-{R}iemannian locally
 symmetric spaces, \href{https://doi.org/10.1016/j.aim.2015.08.029}{\textit{Adv. Math.}} \textbf{287} (2016), 123--236,
 \href{https://arxiv.org/abs/1209.4075}{arXiv:1209.4075}.

\bibitem[15]{Kob89}
Kobayashi T., Proper action on a homogeneous space of reductive type,
 \href{https://doi.org/10.1007/BF01443517}{\textit{Math. Ann.}} \textbf{285} (1989), 249--263.

\bibitem[16]{Kob96}
Kobayashi T., Criterion for proper actions on homogeneous spaces of reductive
 groups, \textit{J.~Lie Theory} \textbf{6} (1996), 147--163.

\bibitem[17]{Kob98}
Kobayashi T., Deformation of compact {C}lifford--{K}lein forms of
 indefinite-{R}iemannian homogeneous manifolds, \href{https://doi.org/10.1007/s002080050153}{\textit{Math. Ann.}}
 \textbf{310} (1998), 395--409.

\bibitem[18]{Ko01}
Kobayashi T., Discontinuous groups for non-{R}iemannian homogeneous spaces, in
 Mathematics Unlimited~-- 2001 and Beyond, \href{https://doi.org/10.1007/978-3-642-56478-9_37}{Springer}, Berlin, 2001, 723--747.

\bibitem[19]{Kob16}
Kobayashi T., Intrinsic sound of anti-de {S}itter manifolds, in Lie Theory and
 its Applications in Physics, \textit{Springer Proc. Math. Stat.}, Vol.~191,
 \href{https://doi.org/10.1007/978-981-10-2636-2_6}{Springer}, Singapore, 2016, 83--99, \href{https://arxiv.org/abs/1609.05986}{arXiv:1609.05986}.

\bibitem[22]{functional-analysis}
Rudin W., Functional analysis, 2nd ed., \textit{International Series in Pure and
 Applied Mathematics}, McGraw-Hill, Inc., New York, 1991.

\bibitem[23]{Sa00}
Salein F., Vari\'et\'es anti-de {S}itter de dimension 3 exotiques, \href{https://doi.org/10.5802/aif.1754}{\textit{Ann.
 Inst. Fourier (Grenoble)}} \textbf{50} (2000), 257--284.

\end{thebibliography}
\end{document}
